% sections/07-casestudy.tex -- Section VII, case study (I4, I3).
% Brief: private/VERA_ICSE2027_writing_brief.md (not versioned), "VII. Case Study". Nothing in this
% section is a result yet. The method is written; the case itself is planned
% (N8: case 25 Sep, run 4 Oct, decision and cost log 9 Oct).
% Honesty: no review may be described as having happened until the record
% exists (handoff item 22). If no real review happens by 2 Oct, D16 switches
% the section to its degraded variant (see the D16 marker below).
% Reviewed 17 Sep (two cached reviews, 29 findings): digest scope, degrees of
% data collection, analysis and member checking, requirement names of Tab. 1.
\psection[1.0]{casestudy}{\armswitch{0.9}{0.9}{1.0}}{I4,I3}{aicc,foundry}{Case Study: Reviewing One Language-Model Use Case}

Section~\ref{sec:results} evaluated VERA under controlled conditions; this section follows one use case through review.
Where an earlier SEIP study described machine-learning engineering across many teams~\cite{amershi2019software}, we follow one review in depth, as a single case with embedded units of analysis~\cite{runeson2009guidelines}.
The case is the review by \aicc{} of one language-model use case\decision{D16}{aicc}{25 Sep}{Which case: one AICC use case, a model from another team's perimeter (not confirmed to be AICC), or three projects; permission to publish the case description and its review record; AICC's agreement to keep an effort log, to a post-decision interview and to check the case description.}, followed as it takes place rather than reconstructed from records\decision{D16}{aicc}{2 Oct}{Real-review go/no-go. If no real review happens by 2 Oct, degrade: a retrospective re-examination of a past AICC decision, retitled Retrospective Case Study: A Past Language-Model Review, labelled as such in the method; cut this section to about 0.6 p; move the freed 0.3 to 0.4 p to VI-D; drop industrial case study from the abstract and reword I4.}.
The units of analysis are the requirements of the language-model specification (Section~\ref{sec:vera}).
The study is descriptive: it records what the evidence covered, what reviewers did with it and what the review cost, without judging the decision itself.
\ifarmwithdrawn
Because the tabular specification was designed but not run (Section~\ref{sec:vera}), this case and the decision study of Section~\ref{sec:decision} carry the paper's evidence from practice, and the case adds a real serving route to the measurability results of Section~\ref{sec:measurability}.
\fi

We collect four sources across the three degrees of data collection that the guidelines distinguish.
First-degree data are an effort log that the \aicc{} reviewers keep during the review and a semi-structured interview with them after the decision, which records how they report reading the evidence\decision{D15}{legal}{25 Sep}{Extend the sign-off and consent text to the case-study interview and the reviewers' effort log.}.
Second-degree data are the VERA run record, with its fallback flags and the content digest of the catalog's version and weights, collected without interacting with reviewers.
Third-degree data are the review record, a work artefact written for the review rather than for this study.
One author maps each finding of the review record to a requirement and an action, a second author checks the mapping, and the \aicc{} reviewers check the case description before publication.

% floats/case.tex -- Tab. 6 (Section VII): evidence and review action per
% requirement in the case study. Hand-built, column width, budget 0.15 p,
% optional (second float to cut, D20).
% Sources, named by the \tbd markers in the paragraph that cites this table:
%   Prior      N7 (AICC coverage before VERA; scope to confirm, D17)
%   VERA       N8 run on the case model, classified with the E4 cause logic
%   Action     N8 review record (publication permission, D16 and D27)
% Rows follow the language-model specification: the nine requirements scored
% per model, then the three scored once on a corpus. Every cell stays \tbdcell
% until its source exists; never type a plausible value.
\registerfloat{tab:case}{table}
\begin{table}[t]
  \centering
  \caption{Case study: evidence and review action per requirement}
  \label{tab:case}
  \tblsetup
  \begin{tabular}{@{}l l c c c@{}}
    \toprule
         & Requirement         & Prior     & VERA      & Review action \\
    \midrule
    \req{01} & Robustness          & \tbdcell & \tbdcell & \tbdcell \\
    \req{02} & Cyber resilience    & \tbdcell & \tbdcell & \tbdcell \\
    \req{06} & Capabilities        & \tbdcell & \tbdcell & \tbdcell \\
    \req{07} & Calibration         & \tbdcell & \tbdcell & \tbdcell \\
    \req{08} & AI disclosure       & \tbdcell & \tbdcell & \tbdcell \\
    \req{09} & Watermark           & \tbdcell & \tbdcell & \tbdcell \\
    \req{10} & Representation bias & \tbdcell & \tbdcell & \tbdcell \\
    \req{11} & Fairness            & \tbdcell & \tbdcell & \tbdcell \\
    \req{12} & Toxicity            & \tbdcell & \tbdcell & \tbdcell \\
    \midrule
    \req{03} & Data adequacy       & \tbdcell & \tbdcell & \tbdcell \\
    \req{04} & Copyright           & \tbdcell & \tbdcell & \tbdcell \\
    \req{05} & Privacy             & \tbdcell & \tbdcell & \tbdcell \\
    \midrule
    \multicolumn{4}{@{}l}{Decision on the use case} & \tbdcell \\
    \bottomrule
  \end{tabular}\\[2pt]
  \parbox{\columnwidth}{\footnotesize\raggedright
    Prior: \cmark{} covered by \aicc{} before VERA, \xmark{} not covered.
    VERA: \stnat{} native, \stfb{} fallback, \stnm{} not measurable, \stna{} not
    applicable; superscript s, h or d: serving, set-up (harness) or by-design cause
    (Section~\ref{sec:measurability}). Action: E escalated, C accepted with
    conditions, A accepted, W waived, \xmark{} no finding. \req{03}--\req{05}:
    scanned on the synthetic corpus, not on the use case's data.}
\end{table}


\parplan[30]{The case is a language-model use case reviewed by \aicc{}; we give its purpose, the model and its serving route, and the reason for its selection.}{Source N8 (aicc; case 25 Sep; where the model runs 28 Sep). One sentence on the use case and the reason it was sent to review. Model and version from a probe (every model version in the vera-foundry config is unknown today). Serving route: Foundry serverless, Azure OpenAI, managed compute, or an AICC environment. The AICC model is not in the vera-foundry deployment list, so say where the run executed and confirm that inputs were public or synthetic. Selection rationale in the terms of runeson2009guidelines and the alternatives set aside under D16; under full and reduced, say why the case is a language-model review and not a Legal/Risk review of a tabular model. No person names; naming the bank and the team depends on D27.}

\parplan[\armswitch{45}{45}{85}]{Table~\ref{tab:case} sets prior \aicc{} coverage\tbd{N7}{aicc}{25 Sep}{AICC coverage before VERA, per requirement; confirm the cyber-resilience and toxicity scope (D17)} against the state the VERA run reached\tbd{N8}{aicc,foundry}{4 Oct}{VERA state per requirement on the case model: native, fallback or not measurable, with cause class; and the corpus used for R03--R05} and the review's action\tbd{N8}{aicc}{9 Oct}{review action per requirement and decision on the use case, from the review record}; each unmeasured requirement is attributed to the serving route, which bounds what an assessor can observe~\cite{casper2024blackbox}, or to a limit of our set-up.}{Sources: N7 and D17 for prior coverage (roadmap claim, unconfirmed: cyber-resilience R02 and toxicity R12 only; scanners and prompt sets derczynski2024garak, gehman2020realtoxicity); N8 run for VERA states; cause classes and the serving-only rule of Section VI-C. Only serving causes support I3; set-up and by-design causes do not, and unknown causes are reported as unattributed and never support I3. R03-R05 are corpus scans on synthetic data: mark them not applicable for the case unless the run scanned a corpus AICC supplied. Do not report as valid: R11 decodingtrust adult, R10 on the HF path and its BBQ fallback, R09 in statistical mode (D8), the R07 heuristic, R05 as a PII recall. No APSEC coverage counts.}

\parplan[\armswitch{40}{40}{60}]{The review record shows which findings were escalated, accepted with conditions or waived, by which role and line of defence, and whether the decision record keeps the run's evidence, as internal-audit practice expects~\cite{raji2020closing,mitchell2019modelcard}.}{Source: N8 review record (aicc, 9 Oct), published only with permission (D16, D27). Counts per action come from Tab. 6; quote at most one short passage. Place each decision in a line of defence once N7 has placed AICC (schuett2025three, iia2020three). Do not describe a production release decision unless the record shows one (handoff item 22). Do not recycle the APSEC reviewer anecdote.}

\parplan[30]{We record whether the choice of fairness criterion arose in the review and who settled it, which links the case to Section~\ref{sec:fairness}.}{Sources: N8 review record; N10 Legal/Risk arbitration (ask 18 Sep, answer 30 Sep); D5. The current LLM fairness probe for R11 is not valid evidence (brief 0.2); if fairness evidence was needed, use the N2 record-classification runs if they exist (D7). If the question did not arise, say so and give the reason stated in the record. No quotation without permission.}

\parplan[\armswitch{30}{30}{50}]{Review cost is reported as reviewer effort and headcount against the reference cost of Section~\ref{sec:context}\decision{D18}{aicc,legal,core}{25 Sep (method), 5 Oct (number)}{Review-cost baseline: logged effort or structured interviews, and its value; the case cost is compared with it.}, kept apart from the compute cost of the run.}{Sources: N8 cost log (aicc, 9 Oct); N9 baseline (aicc, legal; method 25 Sep, number 5 Oct). E5 gives tokens and estimated EUR per run, not review effort: never add the two. Questionnaire Q9 and Q10 are declared effort and must be labelled as such. No time-saved claim without a logged baseline. Neither the APSEC reviewer anecdote nor the APSEC compute figure may appear.}

\parplan[25]{The case closes with any change \aicc{} made to its review procedure afterwards, from the follow-up interview.}{Source: AICC interview after the decision, by 14 Oct at the latest (aicc). Report only changes made or formally scheduled; if none by then, write one sentence saying so and give the freed filler words to the cost paragraph. Tie each change to a lesson in Section VIII. Wording stays introduced into the review process unless the record supports a stronger claim.}
